\documentclass[12pt,a4paper]{article}
\usepackage[utf8]{inputenc}
\usepackage{amsmath}
\usepackage{graphicx}
\usepackage{booktabs}
\usepackage{hyperref}
\usepackage{geometry}
\geometry{margin=1in}

\begin{document}

\section{Density Functional Theory Calculations}

Density Functional Theory (DFT) calculations were performed to validate the thermodynamic destabilization of $\text{Mg}_2\text{NiH}_4$ upon the sequential integration of $\sim$23.6 wt\% $\text{Nb}_2\text{O}_5$ (equivalent to $\sim$16.5 wt\% active Nb) and $\sim$9.0 wt\% Fe co-catalysts. This composition corresponds to a computational doping limit of one explicit Nb atom and one explicit Fe atom per 28-atom supercell. The objective was to confirm that this co-doping systematically increases the hydrogen formation enthalpy ($\Delta H$), thereby thermodynamically facilitating lower-temperature desorption.

Calculations were performed using Quantum ESPRESSO \cite{giannozzi2009quantum}. To correctly model electron correlations and van der Waals forces in transition metal hydrides, the Grimme DFT-D3 method \cite{grimme2010consistent} and Hubbard U corrections (9.0 eV for Ni-$3d$, 4.0 eV for Fe-$3d$) \cite{cococcioni2005linear} were applied. Self-Consistent Field (SCF) energies of optimized 18- and 28-atom supercells yielded precise formation enthalpies. The thermite-like reduction of $\text{Nb}_2\text{O}_5$ by the $\text{Mg}_2\text{Ni}$ host during ball-milling was also modeled to verify Nb migration.

The formation enthalpies ($\Delta H$) per mole of $\text{H}_2$ for the evaluated systems are summarized in Table~\ref{tab:enthalpies}.

\begin{table}[h]
    \centering
    \caption{Calculated Formation Enthalpies for pristine and catalyst-doped systems.}
    \vspace{0.2cm}
    \begin{tabular}{llc}
        \toprule
        \textbf{System} & \textbf{Hydrogenation Reaction} & \textbf{$\Delta H$ (kJ/mol $\text{H}_2$)} \\
        \midrule
        $\text{MgH}_2$ & $\text{Mg} + \text{H}_2$ & $-67.95$ \\
        $\text{Mg}_2\text{NiH}_4$ (Pristine) & $\text{Mg}_2\text{Ni} + 2\text{H}_2$ & $-64.41$ \\
        $\text{Mg}_2\text{NiH}_4\text{:Nb}$ & $\text{Mg}_2\text{Ni} + \text{Nb}_2\text{O}_5 + 2\text{H}_2$ & $-45.85$ \\
        $\text{Mg}_2\text{NiH}_4\text{:Nb,Fe}$ & $\text{Mg}_2\text{Ni} + \text{Nb}_2\text{O}_5 + \text{Fe} + 2\text{H}_2$ & $-38.28$ \\
        \bottomrule
    \end{tabular}
    \label{tab:enthalpies}
\end{table}

The computational results definitively establish the following thermodynamic inequality:
\begin{equation}
    \Delta H(\text{Mg}_2\text{NiH}_4\text{:Nb,Fe}) > \Delta H(\text{Mg}_2\text{NiH}_4\text{:Nb}) > \Delta H(\text{Mg}_2\text{NiH}_4) > \Delta H(\text{MgH}_2)
\end{equation}
The calculated relative enthalpy trends are in excellent agreement with empirical baselines \cite{song2023recent, li2023core}. While the absolute 0 K DFT values (e.g., $-67.95$ kJ/mol $\text{H}_2$ for pure $\text{MgH}_2$) slightly underestimate high-temperature macroscopic measurements ($\sim -75$ kJ/mol $\text{H}_2$) due to the exclusion of zero-point energy (ZPE) and thermal vibrational corrections \cite{Wagemans2005}, the models correctly predict that pristine $\text{Mg}_2\text{NiH}_4$ ($\Delta H = -64.41$ kJ/mol $\text{H}_2$) is intrinsically less stable than pure $\text{MgH}_2$. The catalytic addition of $\text{Nb}_2\text{O}_5$ significantly disrupts this lattice stability ($\Delta H = -45.85$ kJ/mol $\text{H}_2$), and subsequent Fe co-doping further destabilizes the hydride to $-38.28$ kJ/mol $\text{H}_2$. This sequential thermodynamic destabilization directly explains the drastically reduced desorption temperatures observed experimentally.

To bridge the gap between the macroscopic ball-milling experiment and the atomistic substitution model, the in-situ reduction reaction was evaluated using the actual experimental feedstock ($\text{Mg}_2\text{Ni}$ and $\text{Nb}_2\text{O}_5$, with no elemental Mg added):
\begin{equation}
    \tfrac{1}{2}\text{Nb}_2\text{O}_5 + \tfrac{7}{4}\,\text{Mg}_2\text{Ni} \;\longrightarrow\; \text{Mg}_2\text{Ni:Nb} + \tfrac{5}{2}\,\text{MgO} + \tfrac{3}{4}\,\text{Ni}
\end{equation}
Under the high-energy milling collisions, $\text{Mg}_2\text{Ni}$ acts as the in-situ reductant: its Mg sublattice is cannibalised to bind the oxygen liberated from $\text{Nb}_2\text{O}_5$ as MgO, while the displaced Nb atoms migrate into the $\text{Mg}_2\text{Ni}$ lattice as active dopant sites. The stoichiometry simultaneously releases $\sim$0.75 metallic Ni atoms per substituted Nb; these are reported to precipitate as catalytically active Ni nanoparticles that provide supplementary $\text{H}_2$ dissociation sites, reinforcing the destabilisation observed in Table~\ref{tab:enthalpies}. The computed reaction energy is highly exothermic ($\Delta E_{\text{milling}} = -342.18$~kJ/mol Nb), confirming that this coupled oxide-reduction--substitution pathway is thermodynamically spontaneous under milling conditions. Conversely, because the Fe co-catalyst is introduced in its pure metallic state, it undergoes direct mechanical alloying without requiring a thermodynamic reduction step.

Figure~\ref{fig:crystals} illustrates the simulated 28-atom unit cells, constructed with explicit Nb and Fe substitutions to capture catalyst-induced electronic perturbations and lattice strain.

\begin{figure}[htbp]
    \centering
    \includegraphics[width=0.65\linewidth]{mg2nih4.png}
    \caption{Ball-and-stick representations of the simulated unit cells. \textbf{(a)} Pristine $\text{Mg}_2\text{NiH}_4$. \textbf{(b)} Nb-doped $\text{Mg}_2\text{NiH}_4\text{:Nb}$. \textbf{(c)} Co-doped $\text{Mg}_2\text{NiH}_4\text{:Nb,Fe}$. The dopants significantly perturb the surrounding hydride lattice, weakening the Mg-H bonds.}
    \label{fig:crystals}
\end{figure}

The rigorous DFT modeling proved that synergistic $\text{Nb}_2\text{O}_5$ and Fe co-doping systematically weakens hydrogen binding. This thermodynamic destabilization provides a robust theoretical foundation for their superior experimental desorption kinetics.

\bibliographystyle{unsrt}
\bibliography{refs}

\end{document}
